\chapter{Your Toolkit and Environment Setup}
\label{ch:toolkit}

\chapterguide%
  {Create folders and run basic \textbf{terminal commands}.}
  {organize data and notebooks, configure an environment, install dependencies, check that everything works, and run labs reproducibly.}

Every lab in this course runs from the same folder, in the same Python environment. Twenty minutes spent here saves hours of chasing path errors later.

\section{The course project layout}

Your project root is the course folder. In this repository it is called \texttt{MACHINE LEARNING/}; name your own copy without spaces, for example \texttt{UTP-MACHINE-LEARNING/}, because spaces in paths cause avoidable trouble in terminals. Inside it you will find:

\noindent
\begin{minipage}[t]{0.47\linewidth}
\begin{lstlisting}[style=utpterminal]
all_datasets/       one folder per dataset
  README.md         the index
  supervised/       Labs 1-3: the answer is known
    bank_marketing/
    hotel_bookings/ hotel_bookings.zip,
                    README.md
    lending_club/
    sms_spam/
    telco_customer_churn/
    used_cars/
  unsupervised/     Labs 4-5: find structure
    creditcard_fraud_ulb/
    online_retail_ii/
    wood_surface_defects/
  advanced/         Labs 5, 7, 8
    banking77/
    electricity_load/
\end{lstlisting}
\end{minipage}\hfill
\begin{minipage}[t]{0.51\linewidth}
\begin{lstlisting}[style=utpterminal]
all_experiments/           one folder per lab:
  course_setup.py          the setup every lab shares
  lab01_features/          notebook, .py twin, README
  lab02_supervised_workflow/ (+ predict_churn.py)
  lab03_supervised_models/
  lab04_clustering_pca/
  lab05_anomaly_limited_labels/
  lab06_reinforcement/
  lab07_deep_learning_cnn/
  lab08_deep_learning_lstm/
  book_checks/
    demo_study_log.py       Chapters 2,4,9,10,16,17
    demo_foundations.py     Chapters 3,5-9
    demo_model_families.py  Chapters 10-14
    demo_unsupervised.py    Chapters 15-19
    demo_advanced.py        Chapters 20-21
    check_your_turn.py      Your-turn answers
models/   results/lab01/ ... lab08/   requirements.txt
\end{lstlisting}
\end{minipage}

{\footnotesize Every dataset folder holds its file and a README, its data card: source, licence, columns and quirks. Keep the ZIP files zipped: pandas reads the CSV inside each one. Every lab folder holds one worked notebook, a \texttt{.py} twin that runs the same code, and a README naming its data, outputs and run time; each starts with \texttt{course\_setup.py}, which finds the course folder. \texttt{predict\_churn.py} uses the model saved in Lab~2. Notebooks write models and cards to \texttt{models/} and evidence to \texttt{results/}; do not edit generated outputs by hand. The \texttt{book\_checks/} scripts print every worked number in the chapters and check it with an assert.\par}

\section{The datasets}

The course folder ships with all eleven datasets, so there is nothing to download. \texttt{all\_datasets/} groups them by how the course learns from them: \texttt{supervised/} for tables whose answer is known (Labs~1 to~3), \texttt{unsupervised/} for data where the labs look for structure without the answer (Labs~4 and~5), and \texttt{advanced/} for few labels, deep networks and sequences. Every dataset is public and real, and each serves a lab where its quirks teach something. Click a \textbf{dataset name} to read its documentation and licence; each dataset's README records how its file was made from the source.

\begin{utptable}
\footnotesize
\renewcommand{\arraystretch}{1.0}
\setlength{\tabcolsep}{4pt}
\begin{tabularx}{\linewidth}{@{}>{\raggedright\arraybackslash}m{0.40\linewidth}>{\raggedright\arraybackslash}X@{}}
\toprule
\rowcolor{UTPPaleBlue!92}
\utpth{Dataset} & \utpth{Where the labs use it} \\
\midrule
\href{https://huggingface.co/datasets/gemuchu/hotel_bookings}{Hotel Booking Demand} & Lab~1: data card, features, leakage; Lab~2: date-order validation, overbooking threshold \\
\href{https://huggingface.co/datasets/scikit-learn/churn-prediction}{Telco Customer Churn} & Lab~1: repairing types, scaling and neighbors; Lab~2: tuning and a saved model \\
\href{https://huggingface.co/datasets/MLLab-TS/bank_marketing}{Bank Marketing} & Lab~1: prediction timing, feature selection, leakage \\
\href{https://huggingface.co/datasets/KTHESAT/usa-used-car-market-cars-com}{Used Cars (Cars.com)} & Lab~2: linear regression on messy listings; Lab~3: forests and boosting \\
\href{https://huggingface.co/datasets/codesignal/lending-club-loan-accepted}{Lending Club Loans} & Lab~3: KNN, naive Bayes, trees, SVMs \\
\href{https://huggingface.co/datasets/codesignal/sms-spam-collection}{SMS Spam Collection} & Lab~3 Try-it: text classification \\
\href{https://huggingface.co/datasets/attik/Online-Retail-II-UCI}{Online Retail II} & Lab~4: customer segments with $K$-means and DBSCAN \\
\href{https://huggingface.co/datasets/iluvvatar/wood_surface_defects}{Wood Surface Defects} & Lab~4: PCA; Lab~7: dense networks and CNNs \\
\href{https://huggingface.co/datasets/David-Egea/Creditcard-fraud-detection}{Credit Card Fraud (ULB)} & Lab~5: anomaly detection, imbalanced evaluation \\
\href{https://huggingface.co/datasets/PolyAI/banking77}{Banking77} & Lab~5: self-training and active learning \\
\href{https://huggingface.co/datasets/sxj1024/electricity_load_2011_2014}{Electricity Load Diagrams 2011--2014} & Lab~8: LSTM forecasting \\
\bottomrule
\end{tabularx}
\end{utptable}

\subsection{If a file is missing}

\begin{utpsteps}
  \item Copy it again from the course repository into \textbf{its folder} in \texttt{all\_datasets/}, keeping its \textbf{exact name}, including capitalization.
  \item Keep \textbf{ZIP files zipped}: the notebooks read them in place.
  \item Run the notebook from inside the course folder: it finds \texttt{all\_datasets/} by looking upward from where it runs.
\end{utpsteps}

\begin{pitfall}
For \textbf{\texttt{FileNotFoundError}}, check the working directory and filename. Fix the path rather than moving files into the project root.
\end{pitfall}

\section{Two ways to set up: plain Python or Anaconda}

An environment is a private shelf of Python packages for one project. Two projects can then use different versions of the same library without knocking each other over. Route~A uses Python's own \texttt{venv}; Route~B uses conda. Either works, so pick one and stick with it. If you do not already have a preference, take Route~A.

\begin{keyidea}
Use one environment per project; activate it before installing or running packages.
\end{keyidea}

\section{Route A: install Python and configure PATH}

Download Python from the \href{https://www.python.org/downloads/}{\textbf{official downloads page}}.

\begin{description}[leftmargin=1.15in,style=nextline]
  \item[Windows] Run the installer, enable \textbf{\texttt{PATH}}, and reopen the terminal. Use \textbf{\texttt{py}}.
  \item[macOS] Run the \textbf{\texttt{.pkg} installer} and restart Terminal. Keep the \textbf{system Python} unchanged.
  \item[Linux] Install \textbf{Python 3} with your package manager. Do not replace the \textbf{system Python}.
\end{description}

\textbf{Windows PATH fix.} If \texttt{py} is missing, open \textbf{Edit environment variables for your account}, then \textbf{Path > Edit}. Add the installer's command directory, for example:

\begin{lstlisting}[style=utpterminal]
%LocalAppData%\Python\bin
\end{lstlisting}

\textbf{macOS PATH fix.} If the installer placed Python in \texttt{/usr/local/bin} but your terminal cannot find it, add that directory to your Z shell's \textbf{PATH}:

\begin{lstlisting}[style=utpterminal]
echo 'export PATH="/usr/local/bin:$PATH"' >> ~/.zshrc
source ~/.zshrc
\end{lstlisting}

For \textbf{Debian or Ubuntu}:

\begin{lstlisting}[style=utpterminal]
sudo apt update
sudo apt install python3 python3-venv python3-pip
\end{lstlisting}

\textbf{Verify} the installation and executable location:

\begin{lstlisting}[style=utpterminal]
# Windows
py --version
py -m pip --version
where py

# macOS or Linux
python3 --version
python3 -m pip --version
which -a python3
\end{lstlisting}

\section{Route A: create an isolated virtual environment}

\textbf{\texttt{venv}} keeps this project's packages separate from other projects. See the \href{https://docs.python.org/3/library/venv.html}{\textbf{venv documentation}}.

The root-level \textbf{\texttt{requirements.txt}} lists dependencies, including PyTorch (\texttt{torch}) for Labs~7--8. A CPU is sufficient.

\begin{lstlisting}[style=utpterminal]
numpy
pandas
matplotlib
scikit-learn
joblib
jupyter
torch
\end{lstlisting}

\textbf{Create} the environment in the \textbf{project root}:

\begin{lstlisting}[style=utpterminal]
# Windows
py -m venv .venv

# macOS or Linux
python3 -m venv .venv
\end{lstlisting}

\textbf{Activate} it from the same folder:

\begin{lstlisting}[style=utpterminal]
# Windows (PowerShell)
.venv\Scripts\Activate.ps1

# macOS/Linux
source .venv/bin/activate
\end{lstlisting}

\textbf{Install dependencies} after activation:

\begin{lstlisting}[style=utpterminal]
# Windows
py -m pip install --upgrade pip
py -m pip install -r requirements.txt

# macOS or Linux
python3 -m pip install --upgrade pip
python3 -m pip install -r requirements.txt
\end{lstlisting}

\textbf{Install once; activate every session.} Reuse \textbf{\texttt{.venv}} without reinstalling unchanged requirements.

\begin{utpfig}[The environment routine]{The environment routine. Creating \texttt{.venv} and installing the requirements happen once; every later session starts by activating the environment and goes straight back to the notebooks.}{fig:venv-routine}
\begin{tikzpicture}[
  font=\small,
  pystep/.style={draw=UTPBlue, fill=white, rounded corners=2pt, align=center,
               inner sep=4pt, minimum width=3.3cm, minimum height=1.05cm,
               font=\scriptsize\ttfamily\color{UTPNavy}},
  pyonce/.style={pystep, fill=UTPPaleBlue},
  pyevery/.style={pystep, fill=UTPPaleTeal, draw=UTPTeal},
  pycap/.style={font=\scriptsize\bfseries\color{UTPMuted}},
  pynote/.style={font=\scriptsize\color{UTPMuted}, align=center},
  pygo/.style={-{Stealth[length=2.2mm]}, thick, draw=UTPTeal}]

\node[pyonce]  (venv) at (0,0)     {py -m venv .venv};
\node[pyevery] (act)  at (4.05,0)  {activate};
\node[pyonce]  (inst) at (8.10,0)  {pip install -r\\requirements.txt};
\node[pyevery] (work) at (12.15,0) {run notebooks in\\all\_experiments/};

\draw[pygo] (venv) -- (act);
\draw[pygo] (act)  -- (inst);
\draw[pygo] (inst) -- (work);

\node[pycap] at (0,0.92)     {once per project};
\node[pycap] at (4.05,0.92)  {every session};
\node[pycap] at (8.10,0.92)  {once};
\node[pycap] at (12.15,0.92) {every session};

% the return path: a later session rejoins at activation, not at install
\draw[pygo, draw=UTPOrange]
  (work.south) -- ++(0,-0.85) -| (act.south);
\node[pynote, text=UTPOrange] at (8.10,-1.15)
  {next session: \textbf{reactivate}, do not reinstall};

\end{tikzpicture}
\end{utpfig}

\begin{utpfig}[One environment per project]{Separate environments, separate packages. Each project's \texttt{.venv} holds its own copies of the libraries, and the system Python is left alone; activate the intended project environment before running \texttt{pip install}.}{fig:venv-separate}
\begin{tikzpicture}[
  font=\small,
  envbox/.style={draw=UTPBlue, fill=UTPPaleBlue, rounded corners=2pt,
                 align=center, inner sep=4pt, minimum width=4.3cm},
  sysbox/.style={draw=UTPMuted, fill=UTPPaleGray, rounded corners=2pt,
                 align=center, inner sep=4pt, minimum width=4.3cm},
  pkg/.style={font=\scriptsize\ttfamily, align=center},
  lbl/.style={font=\scriptsize\color{UTPMuted}, align=center},
  hit/.style={-{Stealth[length=2.2mm]}, thick, draw=UTPTeal}]

\node[sysbox] (sys) at (0,1.75) {system Python};
\node[pkg,text=UTPMuted,anchor=west] at (2.35,1.75) {leave this alone};

\node[envbox] (a) at (-2.80,0) {\texttt{THIS DIRECTORY/.venv}};
\node[pkg,anchor=north] at (-2.80,-0.52) {numpy\\pandas\\scikit-learn\\jupyter};

\node[envbox] (b) at (2.80,0) {\texttt{some other project/.venv}};
\node[pkg,anchor=north] at (2.80,-0.52) {its own versions\\of the same names};

\draw[hit] (sys) -- (a);
\draw[hit] (sys) -- (b);

\end{tikzpicture}
\end{utpfig}


\section{Route B: Anaconda or Miniconda}

\textbf{Miniconda} supplies Python and \texttt{conda}; \textbf{Anaconda} adds Navigator and preinstalled packages. Download \href{https://docs.anaconda.com/miniconda/}{\textbf{Miniconda}} or \href{https://www.anaconda.com/download}{\textbf{Anaconda}}.

\begin{description}[leftmargin=1.15in,style=nextline]
  \item[Windows] Use default installer options without adding conda to PATH. Open \textbf{Anaconda Prompt} from Start.
  \item[macOS / Linux] Run the \texttt{.pkg} or \texttt{.sh} installer, then close and reopen the terminal. The prompt shows \texttt{(base)}, which means conda is active.
\end{description}

\textbf{Create} one environment for the course, with a Python version the libraries support, then \textbf{activate} it:

\begin{lstlisting}[style=utpterminal]
conda create -n ml python=3.12
conda activate ml
\end{lstlisting}

\textbf{Install dependencies} with the same \texttt{requirements.txt} used in Route A:

\begin{lstlisting}[style=utpterminal]
pip install -r requirements.txt
\end{lstlisting}

Create once, activate each session, and avoid installing into \texttt{(base)}. Check the active environment and Python:

\begin{lstlisting}[style=utpterminal]
conda env list
python -c "import sys; print(sys.executable)"
\end{lstlisting}

\begin{pitfall}
Avoid installing the same package with both conda and pip. Use pip with the course requirements. Recreate a broken environment with \texttt{conda env remove -n ml}, then repeat setup.
\end{pitfall}

\section{Check that everything works}

Run these three commands from the course folder with your environment active. Each one tests something different, so a failure tells you where to look.

\begin{lstlisting}[style=utpterminal]
python -c "import numpy, pandas, matplotlib, sklearn, joblib, torch; print('Packages OK')"
python all_experiments/book_checks/demo_study_log.py
python all_experiments/lab01_features/lab01_features_worked.py
\end{lstlisting}

\begin{utptable}
\begin{tabularx}{\linewidth}{@{}>{\bfseries\leavevmode\color{UTPNavy}}m{0.31\linewidth}m{0.33\linewidth}X@{}}
\toprule
\rowcolor{UTPPaleBlue!92}
\utpth{Command} & \utpth{You should see} & \utpth{It proves} \\
\midrule
The import line & \texttt{Packages OK} & The environment is active and complete \\
\texttt{demo\_study\_log.py} & Last line: \texttt{Every number above matches the book.} & NumPy reproduces the chapters' worked examples \\
\texttt{lab01\_features\_worked.py} & \texttt{raw + engineered validation ROC AUC = 0.855} among the lines & The datasets are in place and scikit-learn works \\
\bottomrule
\end{tabularx}
\end{utptable}

\section{Start a notebook, and read its device line}

Every lab has a worked notebook. Start Jupyter from the course folder with \texttt{jupyter lab}; each notebook finds \texttt{all\_datasets/} by looking upward from its own folder. Open a notebook and choose \textbf{Restart Kernel and Run All Cells}. Most labs take under a minute; Labs~7 and~8 train neural networks, about four minutes each on a laptop CPU. Their Setup cell prints the Python in use and the device it picked: an NVIDIA GPU (\texttt{cuda}), an Apple GPU (\texttt{mps}), or the CPU. To reproduce the printed numbers exactly on a machine with a GPU, hide it before starting Jupyter:

\begin{lstlisting}[style=utpterminal]
# PowerShell
$env:CUDA_VISIBLE_DEVICES = "-1"
# macOS or Linux
export CUDA_VISIBLE_DEVICES=-1
\end{lstlisting}

\begin{troubleshoot}
\begin{description}
  \item[\texttt{Activate.ps1} cannot be loaded] Windows blocks scripts by default. Run \texttt{Set-ExecutionPolicy -Scope CurrentUser RemoteSigned} once in PowerShell, then activate again.
  \item[Conda not found on Windows] Use Anaconda Prompt, or initialize PowerShell there with \texttt{conda init powershell}.
  \item[A notebook cannot import a package] Jupyter is running from a different environment. Activate the course environment, restart Jupyter, and run \texttt{import sys; print(sys.executable)} in a cell to confirm which Python it uses.
  \item[\texttt{Invalid device id} in the Setup cell of Lab~7 or~8] You set \texttt{CUDA\_VISIBLE\_DEVICES} to an empty string. Use \texttt{-1}, as shown above.
  \item[\texttt{OMP: Error \#15}] Two OpenMP runtimes are loaded. The Setup cells of Labs~7 and~8 set \texttt{KMP\_DUPLICATE\_LIB\_OK}; elsewhere, set it before importing PyTorch.
\end{description}
\end{troubleshoot}
